\section{研究背景与目前系统现状}

\subsection{研究背景}

近年来，大语言模型驱动的多智能体（Multi-Agent）系统正逐步从单轮问答走向长链路任务执行。在生命科学场景中，这一趋势尤为明显：系统不再只负责文献检索、知识问答或方案建议，而是开始尝试组织更复杂的科研流程，例如候选生成、结构预测、质量门控、失败恢复以及人工介入等。

对于蛋白设计这类高代价任务而言，系统真正面临的难点并不只是“是否会调用工具”，而在于以下几个方面。首先，任务链条较长，阶段之间存在明显依赖，一个阶段的局部错误往往会影响后续多个阶段。其次，不同阶段的计算代价和失败后果高度不对称，错误决策不仅会带来当前步骤的损失，还可能放大后续恢复成本。再次，运行过程中经常出现信息不充分、候选分歧、局部失败和恢复选择等问题，系统需要不断根据新证据调整执行策略。最后，何时继续自动执行、何时切换到保守路径、何时请求人工介入，本身就是决定系统质量的关键问题。

因此，如果将毕设继续定义为“生命科学助手”或“蛋白设计助手”，那么研究重点就很容易停留在任务分解、工具调用和结果输出这一层。即使在系统内部加入门控、人工介入和恢复机制，课题的主叙事仍然会是“系统如何帮助完成生命科学任务”。若希望体现更明确的计算机专业特征，就需要把问题进一步上升为：\textbf{高代价科学工作流如何被稳定地控制和治理}。

\subsection{当前系统现状}

从现有系统设计来看，它已经不是一个简单的对话式助手，而是一个较完整的多智能体工作流执行系统。

首先，系统采用显式有限状态机（\texttt{FSM}）作为运行核心，要求任务状态显式、状态转移合法且可记录，同时区分外部语义状态和内部执行状态。这意味着系统已经具备较好的执行一致性、可恢复性和可审计性基础。

其次，系统已经形成清晰的角色边界：\texttt{PlannerAgent} 负责生成计划候选，\texttt{ExecutorAgent} 负责执行工具调用和处理恢复链，\texttt{SafetyAgent} 负责风险评估，\texttt{SummarizerAgent} 负责结果汇总。这表明系统已经具备典型多智能体运行时（runtime）结构，而不是简单地将多个提示词串联在一起。

再次，系统已经围绕蛋白设计任务构建了较完整的阶段化流程，目前至少覆盖：
\begin{itemize}
	\item \texttt{S1}：序列候选生成；
	\item \texttt{S2}：结构预测或结构投影；
	\item \texttt{S3}：质量门控；
	\item \texttt{S4}：基于结构条件的细化；
	\item 分层恢复机制：\texttt{retry}、\texttt{patch} 与 \texttt{replan} 构成的恢复链。
\end{itemize}

此外，系统已经具备人工介入（\texttt{HITL}）与恢复链设计，包括若干 \texttt{WAITING\_*} 状态、待确认动作（\texttt{PendingAction}）以及基于快照和事件日志的恢复能力。这说明系统不仅考虑了“如何执行”，也考虑了“如何暂停、恢复和审计”。

最后，系统已经具备一定的实验与评测基础。当前仓库中已有纵向实验框架，可从效果、机制、代价和治理等维度对系统进行分析，并支持增量归因与可复现实验。这一点非常重要，因为它意味着后续研究不必从零构建实验平台。

\subsection{对当前系统的整体评估}

综合来看，当前系统已经具备以下特点：它有较成熟的多智能体运行时骨架，有显式状态机和恢复机制，有蛋白设计这一具体且高代价的工作流场景，也有可继续复用的实验基础。换句话说，它已经具备成为研究平台和实验载体的条件。

但从毕设角度看，这个系统仍然存在一个更关键的问题：\textbf{它更像一个完成度较高的执行平台，而不是一个核心算法问题已经被明确抽象出来的研究课题。}

这一问题主要体现在三个方面。第一，系统虽然已经有打分、风险等级和代价评估等概念，但从实现逻辑上看，它们更多服务于候选排序和展示，仍然偏向启发式聚合，而非一个足够独立、可论证的核心算法。第二，系统当前的主叙事仍然是“生命科学任务输入、多智能体规划执行、工具调用、结果输出”，这意味着即使引入门控、人工介入和恢复链，也仍然容易被看成“生命科学助手”的增强版。第三，系统的许多优点更多体现在工程正确性和执行健壮性上，例如状态一致性、快照恢复、事件审计和人工确认，这些内容虽然重要，却不足以单独形成清晰的算法主线。

\paragraph{小结}
因此，对当前系统的总体判断可以概括为：当前系统已经具备成为研究平台的条件，但尚未形成一个足够清晰、足够独立的核心算法主线。它最大的风险不是“功能还不够多”，而是研究对象还没有从“生命科学助手”提升为更底层、更通用的计算机问题。如果这一点不改变，那么即使继续扩展系统能力，也仍然可能被视为另一生命科学多智能体系统的补充或子集。
